\section{Background and Related Work}
\label{sec:relatedWork}
\subsection{Unstructured observations as process data}
K\"onig et al.~\cite{konig2025unstructured} review how unstructured data enters process mining, including the construction of logs from text, images, and video. Diba et al.~\cite{diba2020extraction} distinguish extraction, correlation, and abstraction. These tasks require decisions about which observations become events, how events are grouped, and which concepts describe them. Andrews et al.~\cite{andrews2020quality} make data quality part of construction in RDB2Log, which helps analysts derive logs from relational data. Together, these perspectives frame log construction as the design of an event representation whose quality must be assessed in relation to its intended use.

% Queue Table 2 for the top of page 4.
\begin{nextpagetopfloat}
\begin{table*}[!t]\begin{adjustwidth}{-\extralength}{0cm}
\caption{Closest construction methods: purpose, input, vocabulary, adaptation, and output. ``Not reported'' identifies information not established by the inspected source, not an absent capability.}
\label{tab:comparison}\footnotesize\setlength{\tabcolsep}{3pt}
\begin{tabularx}{\linewidth}{@{}p{0.14\linewidth}LLLL@{}}\toprule
Method & Purpose and input & Vocabulary origin & Task-specific adaptation & Output \\\midrule
Vid2log~\cite{gavric2024multimodal} & Construction and enrichment from video, with conventional logs & Domain and activity prompts & Pretrained LLaVA; domain prompts; parameter updates not reported & Combined event log \\
Kecht et al.~\cite{kecht2021conversations} & Construction from customer-service conversations & Supplied topics and activities & Pretrained NLI; labeled threshold calibration and hypothesis selection & XES log \\
SnapLog~\cite{pegoraro2026snaplog} & Event discovery from video & Classes supported by labeled examples & Pretrained visual models; classifier adaptation & Events with labels or label distributions \\
Seeth et al.~\cite{seeth2026finetuned} & Construction from textual reports & Allowed activities from source logs & Pretrained LLMs; supervised fine-tuning compared with prompting & XES traces \\
This work & Construction from egocentric factory video & Shared labels derived from recordings & Fixed pretrained models; context and vocabulary induction & Event logs and linked descriptions \\
\bottomrule\end{tabularx}
\end{adjustwidth}\end{table*}
\end{nextpagetopfloat}

\subsection{Video and first-person workflow understanding}
Petersen et al.~\cite{petersen2013workflows} model and track manual workflows from first-person video to support augmented-reality instructions and feedback during task execution. Pham and Xiao~\cite{pham2018workflow} extract assembly steps and operations with semantic descriptions, without pre-labeled training data or pre-trained classifiers. These studies establish first-person workflow understanding as an antecedent. Their goals include following a demonstrated task and recovering its steps.

Retrospective process analysis asks how observations across recordings can support a common description of work. This requires event representations that can be interpreted across tasks and participants. The distinction concerns the analytical objective as well as the camera viewpoint: following a task and analyzing a collection of recorded activities place different demands on the representation.

Video process mining also has several established designs. Knoch et al.~\cite{knoch2020video} use hand trajectories and clustering at a camera-equipped assembly workstation. Kratsch et al.~\cite{kratsch2022video} develop a reference architecture, and Lepsien et al.~\cite{lepsien2023analytics} combine object tracking, action detection, and abstraction. SnapLog~\cite{pegoraro2026snaplog} uses pretrained visual representations, temporal segmentation, and a classifier adapted with labeled examples. Its evaluation includes TUM Kitchen and egocentric EPIC-KITCHENS recordings. Thus, camera placement and supervision must be compared for individual methods. The sensor-data review by Brzychczy et al.~\cite{brzychczy2025sensor} places these efforts within the broader problem of converting observations into process events.

\subsubsection{From video descriptions to process events}
Dense video captioning already combines temporal localization with natural-language descriptions~\cite{krishna2017densecaptioning}. Timestamped narration alone is therefore an established output. Process analysis requires further choices: reusable activity names across recordings, event granularity, selection of operationally relevant episodes, observational case identifiers, source links, and rules for interpreting durations and frequencies.

The unit of a natural-language description need not coincide with the unit required for process analysis. A description can cover several actions, while descriptions of similar actions can use different words. Event construction therefore requires both temporal and semantic abstraction. These choices determine what event counts, durations, and sequences mean.

% Queue Table 3 for the top of page 5.
\begin{nextpagetopfloat}
\begin{table*}[!t]\begin{adjustwidth}{-\extralength}{0cm}
\caption{Sources of temporal information, case relationships, and evaluation references in the compared methods.}
\label{tab:comparison_provenance}\footnotesize\setlength{\tabcolsep}{3pt}
\begin{tabularx}{\linewidth}{@{}p{0.14\linewidth}LLL@{}}\toprule
Method & Origin of temporal information & Origin of case/entity relationships & Evaluation reference \\\midrule
Vid2log~\cite{gavric2024multimodal} & Frame positions and prompted temporal relations & Temporal/semantic matching and resource-based case inference & Assembly case study and log-richness assessment \\
Kecht et al.~\cite{kecht2021conversations} & Recorded message timestamps & Supplied message and conversation identifiers & Labeled messages for classification; process discovery demonstration \\
SnapLog~\cite{pegoraro2026snaplog} & Segmented frame sequences & Video sequences as traces & Video datasets with action and boundary annotations \\
Seeth et al.~\cite{seeth2026finetuned} & Timestamps extracted from generated reports & Report-to-trace pairing from source logs & Source traces and process-level comparisons \\
This work & Inferred intervals relative to recordings & Recording and worker identifiers & Consistency with generated descriptions and structural checks \\
\bottomrule\end{tabularx}
\end{adjustwidth}\end{table*}

\end{nextpagetopfloat}

\subsection{Video-to-log and language-model-based construction}
\label{sec:closest}
Gavric et al.'s Vid2log~\cite{gavric2024multimodal} uses LLaVA to infer a domain and relevant activities from selected frames, obtain textual event information, and combine video-derived data with conventional logs through temporal and semantic matching. It studies log richness and an assembly use case. Both Vid2log and the present work use natural language to connect visual observations to event data. Our focus is on deriving a shared vocabulary from a collection of recordings and preserving the descriptions used to construct the events.

Kecht et al.~\cite{kecht2021conversations} use pretrained natural language inference to assign supplied topics and activities to customer-service messages. They use labeled messages to select classification thresholds and refine hypotheses, without training a new language model. Message identifiers and timestamps are retained. This illustrates an important distinction among construction methods: event names and times may be supplied by the source, selected by an analyst, or inferred from the observations.

Shirali et al.~\cite{shirali2025iot} apply LLMs to physical observations: they abstract IoT readings into high-level events and integrate logs from multiple sensor sources. LLM use for physical-process data therefore predates this study. Work on process-mining assistants~\cite{berti2025benchmark} and database-query generation~\cite{dani2025llm} addresses neighboring tasks with different inputs and outputs.

Seeth et al.~\cite{seeth2026finetuned}, an arXiv preprint, pair generated reports with source traces, supply permitted activities and attributes, and evaluate extracted traces and process models against source-log references. They report benefits from fine-tuning over zero-shot and few-shot prompting in that setting. This provides a comparison in which the vocabulary and evaluation reference come from existing logs. The present work uses fixed pretrained models and derives shared labels from recordings. Differences in supervision and reference data must be considered when comparing the resulting assessments.

Tables~\ref{tab:comparison} and~\ref{tab:comparison_provenance} summarize these differences. They distinguish how a method obtains its vocabulary and adapts its models from how it establishes temporal information, event relationships, and evaluation references.

\subsection{Event abstraction and semantic consistency}
Event abstraction means choosing the level of detail represented by an event. The review by van Zelst et al.~\cite{vanZelst2021abstraction} explains how this choice changes process granularity and compares methods by supervision, input, and uncertainty. Three questions are relevant to construction from video: where an episode begins and ends, which reusable activity classes describe the work, and which class is assigned to each episode. Separating these questions helps distinguish temporal segmentation, vocabulary design, and event classification.

Vocabulary design and event classification can fail in different ways. A relevant concept may be absent from the vocabulary, or an available concept may be assigned to the wrong observation. Further aggregation introduces another interpretive choice by grouping detailed labels into broader categories. Semantic consistency therefore concerns both the meaning assigned to individual events and the meaning retained when those events are analyzed together.

\subsection{Construction, enrichment, and integration}

% Queue Table 4 for the top of page 6.
\begin{nextpagetopfloat}
\begin{table*}[!t]\begin{adjustwidth}{-\extralength}{0cm}
\caption{Neighboring event-data tasks, distinguished by their starting points and intended results. A system can perform more than one task.}
\label{tab:construction_tasks}\footnotesize
\begin{tabularx}{\linewidth}{@{}lLL@{}}\toprule
Task & Starting point & Intended result \\\midrule
Construction~\cite{diba2020extraction} & Observations or documents without a suitable log & An initial event representation \\
Enrichment~\cite{kapugamageeganage2024text2el} & An existing log and additional information & Added or corrected events and attributes \\
Integration~\cite{shirali2025iot,sturm2026multimodal} & Multiple sources describing related behavior & A combined representation with resolved correspondences \\
\bottomrule\end{tabularx}
\end{adjustwidth}\end{table*}
\end{nextpagetopfloat}

Text2EL+~\cite{kapugamageeganage2024text2el} enriches existing logs with events and attributes extracted from text, using semantic and contextual validation and expert guidance. Sturm et al.~\cite{sturm2026multimodal} instantiate parsing, abstraction, extraction, integration, and enrichment using CRM data, email metadata, and call transcripts. Their evaluation involves a sales process and a domain expert. Table~\ref{tab:construction_tasks} separates the tasks these designs can combine.

The present work primarily concerns construction: it derives an initial event representation from recordings. This differs from adding information to an existing log or resolving correspondences among independently recorded sources. Intermediate descriptions can make a construction process easier to inspect, but agreement between representations derived from the same observations does not provide independent corroboration.

\subsection{Quality, uncertainty, and relational representations}
Suriadi et al.~\cite{suriadi2017imperfections} organize recurring event-log imperfections. For logs constructed from unstructured observations, quality depends on both the source evidence and the transformations applied to it. Assessment must therefore consider how information can be omitted, altered, or propagated before it appears in an event record.

Quality annotations and explicitly uncertain events serve different purposes. Goel et al.~\cite{goel2022quality} retain data-quality information for use during analysis. Pegoraro et al.~\cite{pegoraro2021uncertain} represent uncertainty explicitly and compute conformance bounds. SnapLog can retain candidate-label distributions~\cite{pegoraro2026snaplog}. These approaches distinguish documenting known quality issues from representing alternative event interpretations. Consistency checks alone do not establish calibrated probabilities or bounds on downstream error.

Esser and Fahland~\cite{esser2021multidimensional} represent events and entities across multiple dimensions in graph databases. This complements the case-choice concerns of object-centric process mining~\cite{aalst2019objectcentric}. Such representations make the relationships among events explicit, but those relationships still require supporting observations. The range of industrial process-mining goals~\cite{corallo2020industrial} also makes it necessary to align the available evidence with the intended analytical claims.

Building on these methods, we study the construction of shared event representations from egocentric factory recordings. The emphasis is on making the path from observations to process measurements inspectable, so that the resulting descriptions of work can be interpreted in light of their construction.
